\documentclass[12pt, a4paper]{article}
\usepackage{../notesheets}
\usepackage{epsdice}
\usepackage{tikz}
\usepackage{tkz-euclide}
\usetikzlibrary{shapes.geometric, arrows}

\newcommand{\Var}{\operatorname{Var}}
%%%%%%%%%%%%%%%%%%%%%%%%%%%%%%%%%%%%%%%%%%%%%%%%%%
\author{Math 1220}
\title{Notesheet. Section 9.4: Euler's Method} 
\date{}

\begin{document}
\maketitle
\nameline
%%%%%%%%%%%%%%%%%%%%%%%%%%%%%%%%%%%%%%%%%%%%%%%%%%
We cannot solve every \(1\)st order differential equation of the
form \[
  y' = F(x,y)
\]
However, given a specific point, we know the slope of the tangent line
of the solution \(y(x)\) passing through any given point
\((x_0,y_0)\).\\
\begin{tikzpicture}
   \tkzInit[xmax=5,ymax=5,xmin=0,ymin=0]
   \tkzGrid
   \tkzAxeXY
 \end{tikzpicture}\hspace{2in}
 \begin{tikzpicture}
   \tkzInit[xmax=5,ymax=5,xmin=0,ymin=0]
   \tkzGrid
   \tkzAxeXY
 \end{tikzpicture}
 
 Using this information, we can approximate solutions to the IVP \[
   y' = F(x,y), y(x_0) = y_0
 \]
 \begin{ex}
   Approximate the value of \(y(4)\) where \(y(x)\) is
   the solution to the IVP \[
     y'=x-y, y(0) = 1
   \]
 \end{ex}
\pagebreak
 \begin{thrm}[Euler's Method]
   Given a first order \(IVP\) of the form \[
     y' = F(x,y), \ y(x_0) = y_0
   \]
   we can approximate \(y(b)\) in \(n\) steps by
   \begin{enumerate}[label=(\roman*)]
   \item Finding the step size \(h = \frac{b-x_0}{n}\)
   \item Determining the \(x\)-values \(x_i = x_{i-1}+h\) for \(1 \leq
     i \leq n\). \vspace{0.5in}
   \item Determining the \(y\)-values \[
       \begin{cases}
         y_1 = \\
         y_2 = \\
         \vdots\\
         y_i = 
       \end{cases}
     \]
   \end{enumerate}
 \end{thrm}
 \begin{ex}
   Use Euler's method with \(n=1\), \(n=2\), and \(n=4\) to obtain an
   approximation of the solution of the IVP \[
     y' = xy, y(1) = 2
   \]
   when \(x=5\). Which approximation is the most accurate?
 \end{ex}
%%%%%%%%%%%%%%%%%%%%%%%%%%%%%%%%%%%%%%%%%%%%%%%%%% 
\end{document}
